\subsection{Relative density for deterministic presentations}

We use the following deterministic presentation convention in the
relative-density statements.

\begin{definition}[Trim deterministic presentation]
  \label{def:trim-deterministic-presentation}
  A deterministic presentation over a finite alphabet \(\Sigma\) is a
  tuple
  \[
    \mathcal N=(Q_N,\Sigma,\delta_N,q_N,F_N),
  \]
  where \(Q_N\) is finite, \(q_N\in Q_N\), \(F_N\subseteq Q_N\), and
  \(\delta_N:Q_N\times\Sigma\rightharpoonup Q_N\) is a partial
  transition map. A word is accepted if it labels a path from \(q_N\)
  to a state in \(F_N\). Its adjacency matrix is
  \[
    (B_N)_{qq'}
    :=\#\{a\in\Sigma:\delta_N(q,a)=q'\};
  \]
  thus several letters leading from \(q\) to the same state \(q'\) are
  counted with their multiplicity, while unavailable letters contribute
  no edge. The presentation is trim if every state is reachable from
  \(q_N\) and can reach some state of \(F_N\). It is strongly connected
  if the directed graph with edges \(q\to q'\) whenever
  \(\delta_N(q,a)=q'\) for some \(a\in\Sigma\) is strongly connected.

  If \(\mathcal A=(Q_A,\Sigma,\delta_A,q_A,F_A)\) is another
  deterministic automaton over the same alphabet, its synchronous
  product with \(\mathcal N\) has state set \(Q_N\times Q_A\), initial
  state \((q_N,q_A)\), accepting set \(F_N\times F_A\), and the partial
  transition
  \[
    (q,s)\xrightarrow{a}(\delta_N(q,a),\delta_A(s,a))
  \]
  exactly for letters \(a\) for which both displayed transitions are
  defined. Completing \(\mathcal A\) by a sink state, when desired, does
  not change the intersection language.
\end{definition}

We first record the finite-budget statement in the general
multicomponent form. This is the version needed when the deterministic
presentation is trim but its adjacency matrix has several maximal
strongly connected components, possibly connected by transient paths.
The polynomial factors in the statement are exactly the factors absent
from the strongly connected specialization below.

\noindent\textit{\(\lambda\)-active residues.}
Let \(\mathcal N\) be a trim deterministic presentation with adjacency
matrix \(B_N\), let \(\lambda=\rho(B_N)>1\), and let
\(z_m=\abs{L(\mathcal N)\cap\Sigma^m}\). After a period \(D\) has been
fixed, call a residue class \(r\bmod D\) \(\lambda\)-active if
\[
  z_m\not=O(\theta^m)
  \qquad (m\equiv r\pmod D)
\]
for every \(\theta<\lambda\). Equivalently, accepted words in that
residue class have exponential growth rate \(\lambda\). A residue class
can contain infinitely many accepted lengths without being
\(\lambda\)-active: in a reducible presentation those words may be
supplied only by components of smaller spectral radius. The
multicomponent theorem below separates this lower-growth phenomenon
from the Perron-scale comparison.

\begin{theorem}[Multicomponent finite-budget obstruction]
  \label{thm:multicomponent-finite-budget}
  Let \(\mathcal N\) be a trim deterministic presentation,
  not necessarily strongly connected, with adjacency matrix \(B_N\).
  Let \(\lambda:=\rho(B_N)>1\), and put
  \(z_m=\abs{L(\mathcal N)\cap\Sigma^m}\). Then there are an integer
  period \(D\ge1\), a number \(0<\theta<\lambda\), an integer
  \(S\ge0\), and, for every residue \(r\bmod D\), an exponent
  \(s_r\in\{-\infty,0,\dots,S\}\) and a constant \(d_r\ge0\) such that
  \(s_r\ge0\) and \(d_r>0\) exactly on the \(\lambda\)-active residue
  classes, while \(s_r=-\infty\) and \(d_r=0\) on all other residue
  classes. On every residue class
  \[
    z_m
    =
    \lambda^m
    \left(
      d_r m^{s_r}
      +O(m^{s_r-1})
    \right)
    +O(\theta^m m^S),
    \qquad m\equiv r\pmod D .
  \]
  The convention is that the \(O(m^{-1})\) term is absent when
  \(s_r=0\), and the displayed polynomial part is absent when
  \(s_r=-\infty\). Thus non-\(\lambda\)-active residue classes, even
  those with infinitely many accepted lengths, are absorbed into the
  sub-Perron error term.

  For every deterministic automaton \(\mathcal A\) over \(\Sigma\) and
  \[
    a_m(\mathcal A,\mathcal N)
    :=\abs{L(\mathcal A)\cap L(\mathcal N)\cap\Sigma^m},
  \]
  after replacing \(D,S,\theta\) by common choices, there are
  constants \(e_r\ge0\) and exponents
  \(t_r\in\{-\infty,0,\dots,S\}\) such that on every residue class
  \[
    a_m(\mathcal A,\mathcal N)
    =
    \lambda^m
    \left(
      e_r m^{t_r}
      +O(m^{t_r-1})
    \right)
    +O(\theta^m m^S),
  \]
  with the same convention when \(t_r=0\); if the intersection has
  exponential growth \(<\lambda\) on that residue, then \(t_r=-\infty\)
  and the first displayed polynomial part is absent. Moreover, whenever
  \(r\bmod D\) is \(\lambda\)-active, either \(t_r=-\infty\) or
  \(t_r\le s_r\), and if \(t_r=s_r\) then \(0\le e_r\le d_r\).

  Consequently, let \(T_m\subseteq L(\mathcal N)\cap\Sigma^m\) satisfy
  \[
    \alpha\frac{\lambda^m}{m}
    \le \abs{T_m}
    \le \beta\frac{\lambda^m}{m}
  \]
  for all sufficiently large \(m\) in every \(\lambda\)-active residue
  class, with fixed \(\alpha,\beta>0\). For every fixed deterministic automaton
  \(\mathcal A\) there is a constant
  \(C_{\mathcal A,\mathcal N,T}>0\) such that
  \[
    \abs{(L(\mathcal A)\cap L(\mathcal N)\cap\Sigma^m)\triangle T_m}
    \ge
    C_{\mathcal A,\mathcal N,T}\frac{\lambda^m}{m}
  \]
  for all sufficiently large \(m\) in \(\lambda\)-active residue
  classes. On each \(\lambda\)-active residue class either the recall
  against \(T_m\) tends to \(0\) exponentially or polynomially, or else
  the precision is \(O(1/m)\) whenever it is defined.
\end{theorem}

\begin{proof}
  Put \(B_N\) in Frobenius normal form. For a finite non-negative
  integer matrix, the rational generating function
  \(u_N^{\mathsf T}(I-xB_N)^{-1}v_N\) has poles at reciprocals of
  eigenvalues of \(B_N\), and the coefficients are finite sums of
  terms \(p_{\mu}(m)\mu^m\), where \(p_\mu\) is a polynomial whose
  degree is one less than the size of the largest Jordan block attached
  to \(\mu\), plus exponentially smaller terms. Equivalently, this is
  the Frobenius normal form or transfer-matrix expansion for regular
  languages; see \cite{SalomaaSoittola1978,BerstelReutenauer2011}.
  For a non-negative matrix, the eigenvalues of modulus \(\lambda\)
  that contribute to accepted paths are roots of unity times
  \(\lambda\), arising from maximal components and chains between them.
  Taking \(D\) to be a common period of those components makes the
  peripheral factors constant on each residue class. On a fixed residue
  \(r\bmod D\), look only at the terms whose exponential factor has
  modulus \(\lambda\). If no such term contributes to accepted paths in
  that residue, then the whole residue-class count has exponential
  growth \(<\lambda\); after enlarging \(\theta\) below \(\lambda\), it
  is included in \(O(\theta^m m^S)\), and we set \(s_r=-\infty\) and
  \(d_r=0\). This case includes residues with infinitely many accepted
  lengths supplied only by smaller components.

  If a Perron-scale term does contribute, collect the largest remaining
  polynomial degree as \(s_r\). The leading coefficient \(d_r\) is the
  sum of the non-negative leading contributions of maximal
  component-chains realizing that residue at that degree. This sum is
  positive, because at least one such Perron-scale chain contributes
  and no cancellation is possible for counts of paths in a
  non-negative matrix. These are precisely the \(\lambda\)-active
  residues, and the displayed expansion for \(z_m\) follows.

  Form the synchronous product of \(\mathcal N\) with \(\mathcal A\),
  trim unreachable or non-coaccessible states, and let \(B'\) be its
  adjacency matrix. Its accepted-word count is
  \(a_m(\mathcal A,\mathcal N)=u'^{\mathsf T}(B')^m v'\), so the same
  transfer-matrix expansion applies. Its exponential growth rate is at
  most \(\lambda\), because every accepted product word projects
  injectively to a word of \(L(\mathcal N)\). If the product growth rate
  is \(<\lambda\) on a residue class, it is absorbed into
  \(O(\theta^m m^S)\) after increasing \(\theta\) below \(\lambda\) and
  increasing \(S\). If the product has growth \(\lambda\), its leading
  polynomial degree on a \(\lambda\)-active residue cannot exceed \(s_r\);
  otherwise \(a_m(\mathcal A,\mathcal N)\le z_m\) would fail along that
  residue for large \(m\). The same inequality gives \(e_r\le d_r\)
  when \(t_r=s_r\).

  Fix a \(\lambda\)-active residue class and put
  \(A_m=L(\mathcal A)\cap L(\mathcal N)\cap\Sigma^m\). If
  \(t_r=-\infty\), then \(\abs{A_m}=O(\theta^m m^S)\), and hence
  \[
    \frac{\abs{A_m\cap T_m}}{\abs{T_m}}
    \le
    O(m^{S+1}(\theta/\lambda)^m),
  \]
  so recall decays exponentially and
  \(\abs{A_m\triangle T_m}\ge(\alpha/2)\lambda^m/m\) for all large
  \(m\) in the residue.

  Suppose next that the first non-zero leading term in the expansion of
  \(a_m(\mathcal A,\mathcal N)\) on this residue is
  \(\gamma_r\lambda^m m^{u_r}\), with \(\gamma_r>0\) and \(u_r\ge0\).
  Then, for all large \(m\) in the residue,
  \(\abs{A_m}\ge(\gamma_r/2)\lambda^m m^{u_r}\). Since
  \(\abs{T_m}\le\beta\lambda^m/m\),
  \[
    \abs{A_m\triangle T_m}
    \ge \abs{A_m}-\abs{T_m}
    \ge \frac{\gamma_r}{4}\lambda^m m^{u_r}
    \ge \frac{\gamma_r}{4}\frac{\lambda^m}{m},
  \]
  after enlarging the threshold. Whenever \(\abs{A_m}>0\),
  \[
    \frac{\abs{A_m\cap T_m}}{\abs{A_m}}
    \le
    \frac{2\beta}{\gamma_r}\frac{1}{m^{u_r+1}}
    \le
    \frac{2\beta}{\gamma_r}\frac1m .
  \]
  If no such non-zero polynomial term exists, the residue is in the
  exponentially smaller case already handled. There are only finitely
  many \(\lambda\)-active residue classes, so taking the minimum of the
  positive constants obtained above proves the uniform finite-budget
  bound. The theorem makes no prime-scale conclusion on
  non-\(\lambda\)-active residues, where \(L(\mathcal N)\) itself has
  sub-Perron growth and a target of size \(\asymp\lambda^m/m\) cannot
  be a subset of \(L(\mathcal N)\cap\Sigma^m\) for all large \(m\).
\end{proof}

No additional numeration-system formalism is needed for the next
theorem beyond such a trim deterministic presentation whose underlying
directed graph is strongly connected. In the applications below
\(\mathcal N\) is the canonical presentation of a regular numeration
system. Write
\[
  \mathcal N_m:=L(\mathcal N)\cap\Sigma^m,
  \qquad
  z_m:=\abs{\mathcal N_m}.
\]
Let \(B_N\) be the adjacency matrix of \(\mathcal N\), let
\(\lambda:=\rho(B_N)>1\) be its spectral radius, and let \(d_N\ge 1\)
be the period of \(B_N\).

\begin{lemma}[Cyclic-class path threshold]
  \label{lem:cyclic-class-path-threshold}
  Let \(B\) be an irreducible non-negative integer matrix with period
  \(d\) and cyclic decomposition \(C_0,\dots,C_{d-1}\), indexed so that
  \(B\) maps \(C_i\) to \(C_{i+1}\). For each pair of cyclic classes
  \(C_i,C_j\) there is an integer \(N_{ij}\) such that, whenever
  \(n\ge N_{ij}\) and \(n\equiv j-i\pmod d\), every pair
  \(x\in C_i\), \(y\in C_j\) is joined by at least one path of length
  \(n\). The threshold is uniform over the finite pair
  \(C_i\times C_j\).
  This is the elementary cyclic-class input used in the denominator
  positivity argument; it is not a separate contribution.
\end{lemma}

\begin{proof}
  The diagonal block of \(B^d\) on each cyclic class \(C_i\) is
  primitive. Hence for each \(i\) there is an exponent \(E_i\) such that
  \((B^{kd})_{xy}>0\) for every \(x,y\in C_i\) and every \(k\ge E_i\).
  If \(j-i\equiv h\pmod d\), irreducibility and the cyclic
  decomposition imply that for each \(x\in C_i\) and \(y\in C_j\) there
  is a path from \(x\) to \(y\) whose length is congruent to \(h\)
  modulo \(d\). Choose one such path for each finite pair and write its
  length as \(h+dq_{xy}\). Let \(Q_{ij}:=\max_{x,y}q_{xy}\), and let
  \(N_{ij}:=h+d(Q_{ij}+\max_iE_i)\), with
  \(h\in\{0,\dots,d-1\}\).

  For \(n=h+dk\ge N_{ij}\), use the chosen path
  \(x=x_0\to x_1\to\cdots\to x_{h+dq_{xy}}=y\). Its first \(h\) edges
  end at a state \(w\in C_j\). Since \(B^d|_{C_j}\) is primitive, there
  is a path of length \(d(k-q_{xy})\) from \(w\) to \(w\) whenever
  \(k-q_{xy}\ge E_j\). Follow the first \(h\) edges from \(x\) to \(w\),
  insert this return path at \(w\), and then follow the remaining
  \(dq_{xy}\) edges of the chosen path from \(w\) to \(y\). The total
  length is \(h+d(k-q_{xy})+dq_{xy}=n\).
\end{proof}

\begin{lemma}[Periodic growth of trim presentations]
  \label{lem:trim-density}
  With the notation above there exist \(d_N\) non-negative constants
  \(d_0,\dots,d_{d_N-1}\), at least one of which is positive, and
  \(\theta_N\in(0,1)\) such that
  \[
    z_m=\lambda^m\bigl(d_{m\bmod d_N}+O(\theta_N^m)\bigr)
    \qquad (m\to\infty).
  \]
  This is used only for trim strongly connected presentations; no
  multiple-maximal-component extension is asserted here.
\end{lemma}

\begin{proof}
  Since \(\mathcal N\) is trim and strongly connected, \(B_N\) is an
  irreducible non-negative matrix. The periodic Perron--Frobenius
  theorem therefore gives
  \[
    B_N^m
    =\lambda^m\Bigl(\sum_{\ell=0}^{d_N-1}\omega_{d_N}^{\ell m}E_\ell+R_m\Bigr),
    \qquad \|R_m\|\le C\theta_N^m
  \]
  for some spectral projectors \(E_\ell\), constant \(C>0\), and
  \(\theta_N\in(0,1)\), where \(\omega_{d_N}:=e^{2\pi i/d_N}\).
  Equivalently, on cyclic classes \(C_i,C_j\), if
  \(h\equiv j-i\pmod {d_N}\) and \(x\in C_i,y\in C_j\), then the
  irreducible periodic Perron--Frobenius entry asymptotic gives
  \[
    \lambda^{-(kd_N+h)}(B_N^{kd_N+h})_{xy}
      \longrightarrow h_y\nu_x>0
  \]
  for Perron vectors \(h,\nu>0\), normalized on the corresponding cyclic
  classes.
  For \(0\le r<d_N\) set
  \[
    S_r:=\sum_{\ell=0}^{d_N-1}\omega_{d_N}^{\ell r}E_\ell
        =\lim_{k\to\infty}\lambda^{-(kd_N+r)}B_N^{kd_N+r}.
  \]
  Each matrix in the limit is non-negative, so every entry of \(S_r\) is
  non-negative. Grouping the expansion by the residue class of \(m\)
  yields
  \[
    B_N^m=\lambda^m\bigl(S_{m\bmod d_N}+E_m\bigr),
    \qquad \|E_m\|\le C\theta_N^m.
  \]
  Let \(u_N\) and \(v_N\) be the indicator vectors of the initial and
  accepting states, respectively, and define
  \[
    d_r:=u_N^{\mathsf T}S_r v_N\ge 0.
  \]
  The constant \(d_r\) is positive exactly for the residue classes that
  contain arbitrarily long accepted words. Indeed, if
  \(d_r>0\), then some initial-to-accepting entry of \(S_r\) is
  positive; the convergence defining \(S_r\) gives
  \((B_N^{kd_N+r})_{q_N f}>0\) for all sufficiently large \(k\) and
  some \(f\in F_N\), hence accepted words of lengths \(kd_N+r\).
  Conversely, suppose that the residue class \(r\bmod d_N\) contains
  accepted words of arbitrarily large length. Then for infinitely many
  \(k\) there is an accepting state \(f_k\in F_N\) with
  \((B_N^{kd_N+r})_{q_N f_k}>0\). Since \(F_N\) is finite, one
  accepting state \(f\) occurs for infinitely many \(k\). In an
  irreducible non-negative matrix of period \(d_N\), the
  cyclic-class threshold in \Cref{lem:cyclic-class-path-threshold}
  applies, so \(q_N\) and \(f\) lie in cyclic classes whose difference
  is \(r\), and
  \((B_N^{kd_N+r})_{q_N f}>0\) for all sufficiently large \(k\).
  Applying the entry asymptotic displayed above with \(x=q_N\) and
  \(y=f\) gives \((S_r)_{q_N f}>0\), and therefore
  \(d_r=u_N^{\mathsf T}S_rv_N>0\).
  Substituting into the expansion gives
  \[
    z_m=u_N^{\mathsf T}B_N^m v_N
    =\lambda^m\bigl(d_{m\bmod d_N}+O(\theta_N^m)\bigr).
  \]
\end{proof}

For any deterministic automaton \(\mathcal A\) over \(\Sigma\) we write
\[
  a_m(\mathcal A,\mathcal N)
  :=\abs{L(\mathcal A)\cap\mathcal N_m}.
\]

\begin{lemma}[Substochastic product expansion]
  \label{lem:substochastic-product-expansion}
  Let \(\mathcal N\) be a trim strongly connected deterministic
  presentation with adjacency matrix \(B_N\), Perron root
  \(\lambda>1\), and Perron--Frobenius right eigenvector
  \(h\in\RR_{>0}^{Q_N}\). For every deterministic automaton
  \(\mathcal A\) over the same alphabet, complete or partial as in
  \Cref{def:trim-deterministic-presentation}, there exist an integer
  \(D'\ge 1\), non-negative constants \(e_0,\dots,e_{D'-1}\), and
  \(\eta\in(0,1)\) such that
  \[
    a_m(\mathcal A,\mathcal N)
    =\lambda^m\bigl(e_{m\bmod D'}+O(\eta^m)\bigr).
  \]
  This product-expansion lemma supplies the constants used in
  \Cref{thm:parry-density} for fixed trim strongly connected
  presentations. The broader multicomponent finite-budget comparison is
  \Cref{thm:multicomponent-finite-budget}.
\end{lemma}

\begin{proof}
  If necessary, complete \(\mathcal A\) by adjoining a sink state; this
  leaves \(L(\mathcal A)\cap L(\mathcal N)\) unchanged. Form the
  synchronous product automaton of
  \(\mathcal B=\mathcal N\times\mathcal A\) and trim it by deleting
  states that are unreachable from the initial pair or cannot reach an
  accepting pair. If no state remains, then
  \(a_m(\mathcal A,\mathcal N)=0\) for all \(m\), and the assertion
  holds with \(D'=1\) and \(e_0=0\). Otherwise let \(Q'\) be the state
  set of the resulting product \(\mathcal B'\), let \(B'\) be its
  adjacency matrix, and let \(u'\) and \(v'\) be the indicator vectors
  of its initial and accepting states. Trimming does not change the
  intersection counts, so
  \[
    a_m(\mathcal A,\mathcal N)=u'^{\mathsf T}B'^m v'.
  \]

  Let \(\pi\colon Q'\to Q_N\) forget the \(\mathcal A\)-coordinate.
  Define \(\widetilde h_x:=h_{\pi(x)}\), put
  \(\widetilde D:=\operatorname{diag}(\widetilde h)\), and set
  \[
    P':=\lambda^{-1}\widetilde D^{-1}B'\widetilde D.
  \]
  For a state \(x=(q,s)\in Q'\),
  \[
    \sum_{y\in Q'}P'_{xy}
    =\frac{1}{\lambda h_q}\sum_{y\in Q'}B'_{xy}h_{\pi(y)}.
  \]
  Because the product is deterministic letter by letter, each letter
  available from \(q\) in \(\mathcal N\) contributes to at most one
  outgoing edge from \(x\) in the product automaton, and trimming can
  only delete such edges. Hence
  \[
    \sum_{y\in Q'}B'_{xy}h_{\pi(y)}
    \le \sum_{\substack{q\xrightarrow{a}q'\\ \text{in }\mathcal N}} h_{q'}
    =(B_Nh)_q
    =\lambda h_q.
  \]
  Thus \(P'\) is substochastic.

  Adjoin a cemetery state \(\partial\) and define
  \[
    \widehat P:=
    \begin{pmatrix}
      P' & \mathbf{1}-P'\mathbf{1}\\
      0 & 1
    \end{pmatrix}.
  \]
  The matrix \(\widehat P\) is stochastic, and
  \((\widehat P^m)_{xy}=(P'^m)_{xy}\) for all \(x,y\in Q'\). Apply the
  finite Markov-chain canonical decomposition to \(\widehat P\). After
  permuting states, write
  \[
    \widehat P=
    \begin{pmatrix}
      T & R_1 & \cdots & R_s\\
      0 & C_1 & 0 & 0\\
      \vdots & 0 & \ddots & 0\\
      0 & 0 & 0 & C_s
    \end{pmatrix},
  \]
  where \(T\) is the transient block, \(\rho(T)<1\), and each \(C_j\) is
  an irreducible stochastic recurrent class with period \(d_j\). For
  \(m\ge1\),
  \[
    \widehat P^m=
    \begin{pmatrix}
      T^m &
      \sum_{\ell=0}^{m-1}T^\ell R_1C_1^{m-1-\ell} &
      \cdots &
      \sum_{\ell=0}^{m-1}T^\ell R_sC_s^{m-1-\ell}\\
      0 & C_1^m & 0 & 0\\
      \vdots & 0 & \ddots & 0\\
      0 & 0 & 0 & C_s^m
    \end{pmatrix}.
  \]
  The powers \(T^\ell\) decay exponentially, while each \(C_j^m\) has
  the usual periodic Perron--Frobenius limits. Hence, with
  \(D'=\operatorname{lcm}(d_1,\dots,d_s)\), the entrywise limits
  \[
    S_r:=\lim_{k\to\infty}P'^{kD'+r}
    \qquad(0\le r<D')
  \]
  exist on the \(Q'\times Q'\) block. These limits include the
  non-decaying absorption terms from transient states into recurrent
  classes. Since convergence of \(T^\ell\) is exponential and the
  non-peripheral parts of the recurrent blocks decay exponentially,
  there are a constant \(C>0\) and \(\eta\in(0,1)\) such that
  \[
    P'^m=S_{m\bmod D'}+E_m,
    \qquad
    \|E_m\|\le C\eta^m.
  \]
  The residue matrices are non-negative because they are entrywise
  limits of non-negative matrices. This also covers the case in which
  the restricted block is purely transient, where every \(S_r\) is the
  zero matrix.

  Since \(B'^m=\lambda^m\widetilde DP'^m\widetilde D^{-1}\), we obtain
  \[
    a_m(\mathcal A,\mathcal N)
    =\lambda^m\widetilde u^{\mathsf T}P'^m\widetilde v,
    \qquad
    \widetilde u:=\widetilde D u',
    \quad
    \widetilde v:=\widetilde D^{-1}v'.
  \]
  Therefore
  \[
    a_m(\mathcal A,\mathcal N)
    =\lambda^m\bigl(e_{m\bmod D'}+O(\eta^m)\bigr),
    \qquad
    e_r:=\widetilde u^{\mathsf T}S_r\widetilde v\ge 0,
  \]
  as required.
\end{proof}

\begin{theorem}[Background relative density for a trim strongly connected presentation]
  \label{thm:parry-density}
  Let \(\mathcal N\) be a trim strongly connected deterministic
  presentation in the sense of
  \Cref{def:trim-deterministic-presentation}, with Perron root
  \(\lambda>1\), and let
  \(\mathcal A\) be a deterministic automaton over \(\Sigma\). Then
  there exist an integer \(D_N\ge 1\), non-negative constants
  \(c_0,\dots,c_{D_N-1}\), and \(\theta\in(0,1)\) such that for all
  sufficiently large \(m\) with \(z_m>0\),
  \[
    \frac{a_m(\mathcal A,\mathcal N)}{z_m}
    =c_{m\bmod D_N}+O(\theta^m).
  \]
  In particular, either the ratio
  \(a_m(\mathcal A,\mathcal N)/z_m\) decays exponentially fast, or else
  it stays bounded below by a positive constant along an arithmetic
  progression of lengths that contains infinitely many indices with
  \(z_m>0\).
  This theorem is a background conditional-density tool for the stated
  trim strongly connected deterministic setting; its finite-budget use is
  isolated in
  \Cref{cor:relative-finite-budget-obstruction}.
\end{theorem}

\begin{proof}
  This is the standard conditional-density conclusion for deterministic
  presentations, specialized to our notation; compare
  \cite{HanselPerrin1989,Kozik2005Conditional,Koga2019Density,Montoya2025Relative,BertheGouletPerrin2025}.
  The proof below is included only to identify the constants used later:
  \Cref{lem:substochastic-product-expansion} supplies the product
  expansion, and \Cref{lem:trim-density} supplies the positive
  denominator constants on relevant residue classes.
  Let \(h\in\RR_{>0}^{Q_N}\) be a Perron--Frobenius right eigenvector of
  \(B_N\). By \Cref{lem:substochastic-product-expansion}, there are
  \(D'\ge 1\), non-negative constants \(e_0,\dots,e_{D'-1}\), and
  \(\eta\in(0,1)\) such that
  \[
    a_m(\mathcal A,\mathcal N)
    =\lambda^m\bigl(e_{m\bmod D'}+O(\eta^m)\bigr).
  \]

  Now combine this with \Cref{lem:trim-density}. After replacing \(D'\)
  by a common multiple \(D_N\) of \(D'\) and \(d_N\), and regrouping the
  residue classes, there exist non-negative constants
  \(e_0,\dots,e_{D_N-1}\), \(d_0,\dots,d_{D_N-1}\), and \(\theta\in(0,1)\)
  such that
  \[
    a_m(\mathcal A,\mathcal N)
    =\lambda^m\bigl(e_{m\bmod D_N}+O(\theta^m)\bigr),
    \qquad
    z_m=\lambda^m\bigl(d_{m\bmod D_N}+O(\theta^m)\bigr).
  \]
  If \(z_m>0\) for infinitely many \(m\equiv r\pmod {D_N}\), then
  \Cref{lem:trim-density} gives \(d_r>0\). Along every such residue
  class,
  \[
    \frac{a_m(\mathcal A,\mathcal N)}{z_m}
    =\frac{e_r+O(\theta^m)}{d_r+O(\theta^m)}
    =\frac{e_r}{d_r}+O(\theta^m).
  \]
  Setting \(c_r:=e_r/d_r\) on residue classes with \(d_r>0\) proves the
  stated expansion for all sufficiently large \(m\) with \(z_m>0\).

  Finally, if every relevant \(e_r\) vanishes, then
  \(a_m(\mathcal A,\mathcal N)=O((\lambda\theta)^m)\) on every residue
  class with \(z_m>0\), so the ratio decays exponentially. Otherwise
  some relevant \(e_r\) is positive, and along the corresponding
  arithmetic progression the ratio stays bounded below by a positive
  constant.
\end{proof}

Only the following finite-budget consequence is used in the prime-slice
applications.
It is the exact form of the background relative-density theorem needed
below, after the substochastic product expansion and the positive
periodic denominator constants have already been supplied in
\Cref{lem:substochastic-product-expansion,lem:trim-density}.

\begin{corollary}[Finite-budget lemma for a trim presentation]
  \label{cor:relative-finite-budget-obstruction}
  Let \(\mathcal N\) and \(\mathcal A\) satisfy the hypotheses of
  \Cref{thm:parry-density}. For each \(m\) with \(z_m>0\), let
  \(T_m\subseteq\mathcal N_m\). Suppose that there are constants
  \(\alpha,\beta>0\) such that, for all sufficiently large \(m\) with
  \(z_m>0\),
  \[
    \alpha\,\frac{\lambda^m}{m}
    \le \abs{T_m}
    \le \beta\,\frac{\lambda^m}{m}.
  \]
  Put \(A_m:=L(\mathcal A)\cap\mathcal N_m\). Then there is a constant
  \(C_{\mathcal A,\mathcal N,T}>0\) such that
  \[
    \abs{A_m\triangle T_m}
    \ge C_{\mathcal A,\mathcal N,T}\,\frac{\lambda^m}{m}
  \]
  for all sufficiently large \(m\) with \(z_m>0\).

  More precisely, after passing to the residue classes supplied by
  \Cref{thm:parry-density}, every relevant residue class has one of the
  following two behaviours. If its limiting relative-density constant
  is zero, then
  \[
    \frac{\abs{A_m\cap T_m}}{\abs{T_m}}
    \le \frac{\abs{A_m}}{\abs{T_m}}
  \]
  decays exponentially along that residue class. If its limiting
  relative-density constant is positive, then along that residue class
  \[
    \abs{A_m\triangle T_m}\gg \lambda^m,
    \qquad
    \frac{\abs{A_m\cap T_m}}{\abs{A_m}}=O(1/m)
  \]
  whenever the second ratio is defined.
\end{corollary}

\begin{proof}
  Let \(D_N\), \(c_0,\dots,c_{D_N-1}\), and \(\theta\in(0,1)\) be supplied
  by \Cref{thm:parry-density}. Enlarge \(D_N\), if necessary, so that it
  is also a multiple of the period in \Cref{lem:trim-density}. Then for
  every residue class \(r\bmod D_N\) containing infinitely many lengths
  with \(z_m>0\), there is a constant \(d_r>0\) such that
  \[
    z_m=\lambda^m(d_r+O(\theta^m))
    \qquad (m\equiv r\bmod {D_N}),
  \]
  after increasing \(\theta\) within \((0,1)\) if needed. Hence
  \(z_m\asymp\lambda^m\)
  on each relevant residue class.

  Fix such a residue class. If \(c_r=0\), then
  \(\abs{A_m}/z_m=O(\theta^m)\), and therefore
  \[
    \abs{A_m}=O(\lambda^m\theta^m).
  \]
  Choosing \(\eta\) with \(\theta<\eta<1\), the elementary bound
  \(m\theta^m=O(\eta^m)\) gives
  \[
    \frac{\abs{A_m\cap T_m}}{\abs{T_m}}
    \le \frac{\abs{A_m}}{\abs{T_m}}
    =O(m\theta^m)
    =O(\eta^m),
  \]
  so recall against \(T_m\) decays exponentially on this residue class.
  Also, for all sufficiently large \(m\) in this class,
  \[
    \abs{A_m\triangle T_m}
    \ge \abs{T_m}-\abs{A_m}
    \ge \frac{\alpha}{2}\,\frac{\lambda^m}{m}.
  \]

  If \(c_r>0\), then \(\abs{A_m}/z_m\ge c_r/2\) for all sufficiently
  large \(m\equiv r\bmod {D_N}\) with \(z_m>0\). Since
  \(z_m\asymp\lambda^m\) on the residue class, there is a constant
  \(\gamma_r>0\) such that \(\abs{A_m}\ge\gamma_r\lambda^m\). The upper
  bound on \(\abs{T_m}\) then gives, for all sufficiently large \(m\) in
  the class,
  \[
    \abs{A_m\triangle T_m}
    \ge \abs{A_m}-\abs{T_m}
    \ge \frac{\gamma_r}{2}\lambda^m,
  \]
  and, whenever \(\abs{A_m}>0\),
  \[
    \frac{\abs{A_m\cap T_m}}{\abs{A_m}}
    \le \frac{\abs{T_m}}{\abs{A_m}}
    \le \frac{\beta}{\gamma_r}\,\frac{1}{m}.
  \]

  There are only finitely many relevant residue classes. Taking the
  minimum of the positive constants obtained above proves the uniform
  lower bound \(C_{\mathcal A,\mathcal N,T}\lambda^m/m\) for all
  sufficiently large \(m\) with \(z_m>0\), and the two residue-class
  alternatives prove the remaining assertions.
\end{proof}

\noindent
This finite-budget lemma is the only relative-density consequence used
below: the base-\(b\) and Zeckendorf prime slices are exactly the
instances in which \(T_m\) has size comparable to \(\lambda^m/m\) by
\Cref{lem:base-prime-count,lem:zeck-prime-count}.

\begin{remark}
  When \(\mathcal N\) is the trim strongly connected presentation of a
  Parry numeration system or of the Zeckendorf language, the same
  presentation-level statement applies with \(\lambda\) equal to the
  dominant root of the corresponding recurrence. Presentations with
  multiple maximal components are covered by
  \Cref{thm:multicomponent-finite-budget}, where polynomial factors are
  retained. Nondeterministic or weighted presentations are not treated
  here; see \cite{CharlierRampersadRigoWaxweiler2010} for examples where
  multicomponent behaviour genuinely occurs.
\end{remark}

\begin{remark}[Prime automaticity]
  Shen's prime-automata nonexistence theorem
  \cite{Shen2022} gives a qualitative finite-state obstruction for primes
  in ordinary positional notation. Dubbe \cite[Theorem~1]{Dubbe2024}
  varies the automaton with the cutoff: if \(A(x)\) is the minimum number
  of states needed for exact prime recognition up to \(x\), then
  \[
    A(x)\ge
    x\exp\!\bigl(-c(\log\log x)^2\log\log\log x\bigr)
  \]
  for all sufficiently large \(x\). Our quantifiers are different. We
  fix one DFA first and obtain lengthwise symmetric-difference and
  recall/precision lower bounds on arithmetic progressions of lengths;
  the Zeckendorf section gives the analogous fixed-DFA obstruction in that
  syntax. These results do not improve Dubbe's growing-state lower bound,
  while the cited ordinary-base theorems do not state our fixed-DFA
  residue-class error bounds or their Zeckendorf analogue.
\end{remark}

\begin{remark}[Relation to conditional-density literature]
  Classical treatments of densities of rational languages already appear
  in Hansel and Perrin~\cite{HanselPerrin1989} and in the conditional
  and relative-density analyses of Kozik~\cite{Kozik2005Conditional},
  Koga \cite{Koga2019Density}, Montoya \cite{Montoya2025Relative}, and
  Berth\'e--Goulet-Ouellet--Perrin \cite{BertheGouletPerrin2025}.
  Thus \Cref{thm:dfa-density-dichotomy,thm:parry-density} are used here
  as standard inputs for the fixed-DFA and strongly connected
  deterministic settings. The multicomponent theorem keeps the
  polynomial factors that those inputs suppress. The contribution is
  the conversion of these residue-class expansions into the
  finite-budget prime obstructions recorded in
  \Cref{thm:multicomponent-finite-budget,cor:relative-finite-budget-obstruction,cor:dfa-prime-symmetric-diff,cor:dfa-prime-recall-precision}.
\end{remark}
